% Zumdahl Chapter 13 -- Chemical Equilibrium. ELABORATED notes, 5 blocks.
\documentclass{apchem-notes}

\apcunitword{Chapter}
\apcunit{13}
\apcunittitle{Chemical Equilibrium}
\apcdoctitle{Guided Notes}
\apcsubtitle{Zumdahl \S13.1--13.7 \textbullet\ PDF pp.~650--695 \textbullet\ 5 blocks}

\begin{document}
\apctitleblock

\begin{apcnote}[How this chapter fits the AP course]
Chapter 13 covers \textbf{CED 7.1--7.10} in the textbook's own order --- about
80\% of Unit 7, with nothing off-syllabus.

The two topics it does \emph{not} contain are \textbf{7.11 solubility
equilibria ($K_{sp}$)} and \textbf{7.12 the common-ion effect}, which live in
Zumdahl \S16.1 and \S15.1. Plan to pick those up separately.

\textbf{One connection worth making explicitly:} the ICE tables in Block 4
are the BCA tables from Chapter 3 with new labels. Before/Change/After
becomes Initial/Change/Equilibrium, and the change row is still driven by
coefficient ratios. The only new idea is that the reaction stops partway
instead of running to completion.
\end{apcnote}

% =====================================================================
\blocklesson{The Equilibrium Condition}{Zumdahl \S13.1}

\begin{objectives}
  \item Describe chemical equilibrium as a dynamic condition.
  \item Explain why concentrations become constant but not equal.
\end{objectives}

\reading{Zumdahl \S13.1}{651--654}

\begin{warmup}[3]
  \item Rate law for the elementary step \ce{A + B -> C}:
        \answerblank[3.4cm]{$k[\ce{A}][\ce{B}]$}
  \item Does a catalyst change $\Delta H$? \answerblank[1.6cm]{no}
  \item Balance: \ce{N2 + H2 -> NH3}:
        \answerblank[4.4cm]{\ce{N2 + 3 H2 -> 2 NH3}}
\end{warmup}

\segment{INSTRUCTION A}{25}{What equilibrium actually is}

\notesectionz{Dynamic, not static}{13.1}
\practice{1}

Start with pure reactants. The forward reaction begins fast and slows as
reactants are consumed. Products accumulate, so the reverse reaction speeds
up. Eventually the two rates become
\answerblank[1.6cm]{equal}, and from that moment the concentrations stop
changing.

\[
  \text{at equilibrium:}\qquad
  \text{rate}_{\text{forward}} = \answerblank[3cm]{$\text{rate}_{\text{reverse}}$}
\]

\begin{center}
\begin{tikzpicture}
\begin{axis}[width=9.6cm, height=4.6cm,
    xlabel={time}, ylabel={rate},
    xtick=\empty, ytick=\empty, xmin=0, xmax=10, ymin=0, ymax=10,
    axis lines=left, legend style={draw=none, font=\sffamily\scriptsize,
      at={(0.98,0.75)}, anchor=east}]
  \addplot[seriesA, samples=100, domain=0:10] {3.5+5*exp(-0.55*x)};
  \addlegendentry{forward}
  \addplot[seriesB, samples=100, domain=0:10] {3.5-3.5*exp(-0.55*x)};
  \addlegendentry{reverse}
  \draw[densely dotted] (axis cs:0,3.5) -- (axis cs:10,3.5);
\end{axis}
\end{tikzpicture}
\end{center}

\begin{apctrap}
Equilibrium is \term{dynamic}, not static. Both reactions continue at full
speed --- they simply cancel out. Nothing has stopped; the traffic in both
directions is merely balanced.

And ``constant'' does not mean ``equal.'' Concentrations of reactants and
products settle at values that are typically very different from one
another. Saying ``at equilibrium the concentrations are equal'' is the
single most common error in this unit.
\end{apctrap}

\segment{GUIDED PRACTICE}{15}{True or false, with reasons}

\begin{enumerate}[leftmargin=*,itemsep=0.35em]
  \item At equilibrium, the forward reaction stops.
        \answerblank[7.5cm]{False --- both continue at equal rates}
  \item At equilibrium, concentrations are constant.
        \answerblank[5.5cm]{True}
  \item At equilibrium, $[\text{reactants}] = [\text{products}]$.
        \answerblank[5.5cm]{False --- constant, not equal}
  \item Equilibrium can be reached from either direction.
        \answerblank[5.5cm]{True}
\end{enumerate}

\segment{INSTRUCTION B}{20}{Approaching from either side}

\notesectionz{The same destination}{13.1}
\practice{4}

A crucial property: for a given temperature, starting with pure reactants or
with pure products leads to the \answerblank[2.6cm]{same} equilibrium
condition. Only the \emph{route} differs.

\begin{apcnote}
This is why equilibrium is described by a constant rather than a
procedure. The system does not remember how it got there --- which is the
same path-independence you met with state functions in Chapter 6.
\end{apcnote}

\segment{APPLICATION}{20}{Reasoning about the approach}

\begin{enumerate}[leftmargin=*,itemsep=0.4em]
  \item Sketch how $[\ce{N2}]$, $[\ce{H2}]$, and $[\ce{NH3}]$ change over
        time starting from pure \ce{N2} and \ce{H2}, and mark where
        equilibrium is reached. \workspace[3.6cm]
        \keyonly{\begin{apcanswer}\ce{N2} and \ce{H2} fall (H2 three times as
        steeply), \ce{NH3} rises from zero; all three level off and stay
        constant thereafter at unequal values.\end{apcanswer}}
  \item A sealed flask of \ce{N2O4} is left until no further colour change
        occurs, yet a chemist insists both reactions are still running.
        How could that be demonstrated? \answerlines[3]{Introduce isotopically
        labelled molecules and they will appear on both sides over time,
        showing that interconversion continues even though the net
        concentrations no longer change.}
\end{enumerate}

\begin{exitticket}
Explain the difference between saying concentrations are ``constant'' and
saying they are ``equal'' at equilibrium.
\answerlines[2]{Constant means unchanging over time, which is true.
Equal would mean the same numerical value, which is generally false --- the
equilibrium values depend on $K$ and the starting amounts.}
\end{exitticket}

% =====================================================================
\blocklesson{The Equilibrium Constant}{Zumdahl \S13.2--13.3}

\begin{objectives}
  \item Write $K$ expressions from balanced equations.
  \item Convert between $K_c$ and $K_p$.
\end{objectives}

\reading{Zumdahl \S13.2--13.3}{654--660}

\begin{warmup}[3]
  \item At equilibrium the two rates are:
        \answerblank[2cm]{equal}
  \item Is equilibrium static or dynamic? \answerblank[2cm]{dynamic}
  \item $R$ in \si{\liter\atm\per\mole\per\kelvin}:
        \answerblank[2.4cm]{0.08206}
\end{warmup}

\segment{INSTRUCTION A}{25}{The law of mass action}

\notesectionz{Writing $K$}{13.2}
\practice{5}

For the general reaction $j\ce{A} + k\ce{B} \rightleftharpoons l\ce{C} + m\ce{D}$:
\[
  K = \answerblank[4.6cm]{$\dfrac{[\ce{C}]^l[\ce{D}]^m}{[\ce{A}]^j[\ce{B}]^k}$}
\]

\answerblank[2cm]{Products} over reactants, each raised to its
\answerblank[2cm]{coefficient}.

\begin{apcnote}
Here the coefficients \emph{do} become exponents --- unlike rate laws in
Chapter 12, where they never do. The difference is that $K$ describes a
thermodynamic end state determined by the overall equation, while a rate law
describes a mechanism. Two different questions, two different rules.
\end{apcnote}

\notesub{Manipulating $K$}

\renewcommand{\arraystretch}{1.35}
\noindent\begin{tabular}{@{}p{5.4cm}p{7.4cm}@{}}
\toprule
\textbf{If you\ldots} & \textbf{then the new constant is\ldots}\\
\midrule
reverse the reaction & \answerblank[1.8cm]{$1/K$}\\
multiply coefficients by $n$ & \answerblank[1.6cm]{$K^n$}\\
add two reactions & \answerblank[2.6cm]{$K_1 \times K_2$}\\
\bottomrule
\end{tabular}

\segment{GUIDED PRACTICE}{15}{Write the expressions}

\begin{enumerate}[leftmargin=*,itemsep=0.35em]
  \item \ce{N2(g) + 3 H2(g) <=> 2 NH3(g)}:
        \answerblank[5.4cm]{$[\ce{NH3}]^2/([\ce{N2}][\ce{H2}]^3)$}
  \item \ce{2 SO2(g) + O2(g) <=> 2 SO3(g)}:
        \answerblank[5.4cm]{$[\ce{SO3}]^2/([\ce{SO2}]^2[\ce{O2}])$}
  \item If $K = 4.0$ for a reaction, $K$ for the reverse is:
        \answerblank[2cm]{0.25}
\end{enumerate}

\segment{INSTRUCTION B}{20}{$K_p$ and $K_c$}

\notesectionz{Equilibrium in terms of pressure}{13.3}
\practice{5}

For gases it is often easier to measure pressure than concentration, so
$K_p$ uses partial pressures in place of molar concentrations.

Since $C = P/RT$ for an ideal gas, substituting gives
\[
  K_p = \answerblank[3cm]{$K_c(RT)^{\Delta n}$}
  \qquad \Delta n = \answerblank[8.3cm]{(mol gas products) $-$ (mol gas reactants)}
\]

\begin{workedexample}[Worked example 1: when the two are equal]
For \ce{H2(g) + F2(g) <=> 2 HF(g)}: two moles of gas on each side, so
$\Delta n = 0$ and $(RT)^0 = 1$. Therefore
\[ K_p = K_c \]
Zumdahl works the substitution explicitly --- every $RT$ term cancels.
Whenever the gas-mole totals match, the two constants are identical.
\end{workedexample}

\begin{workedexample}[Worked example 2: when they differ]
For \ce{N2O4(g) <=> 2 NO2(g)} at \SI{25}{\celsius}, $K_c =
4.6\times10^{-3}$. Here $\Delta n = 2 - 1 = 1$:
\[
  K_p = K_c(RT)^1 = (4.6\times10^{-3})(0.08206)(298) = 0.11
\]
The two differ by a factor of $RT \approx 24.5$ --- large enough that
quoting the wrong one is a serious error.
\end{workedexample}

\begin{apctrap}
$\Delta n$ counts \emph{gases only}. Solids and liquids are excluded, for
the same reason they are excluded from the $K$ expression itself
(Block 3). And equilibrium constants are conventionally written without
units --- do not attach any.
\end{apctrap}

\segment{APPLICATION}{20}{Conversions}

\begin{enumerate}[leftmargin=*,itemsep=0.4em]
  \item For \ce{N2(g) + 3 H2(g) <=> 2 NH3(g)}, find $\Delta n$ and write the
        relationship between $K_p$ and $K_c$.
        \answerlines[2]{$\Delta n = 2 - 4 = -2$, so $K_p = K_c(RT)^{-2}$,
        equivalently $K_c = K_p(RT)^2$.}
  \item For \ce{2 SO2(g) + O2(g) <=> 2 SO3(g)} at \SI{700}{\kelvin},
        $K_c = 4.3\times10^{6}$. Calculate $K_p$. \workspace[2.2cm]
        \keyonly{\begin{apcanswer}$\Delta n = 2-3 = -1$;
        $K_p = (4.3\times10^{6})(0.08206 \times 700)^{-1} =
        4.3\times10^{6}/57.4 = 7.5\times10^{4}$\end{apcanswer}}
\end{enumerate}

\begin{exitticket}
For which type of reaction does $K_p = K_c$? Give an example.
\answerlines[2]{One where the total moles of gas are the same on both sides
($\Delta n = 0$), such as \ce{H2(g) + I2(g) <=> 2 HI(g)}.}
\end{exitticket}

% =====================================================================
\blocklesson{Heterogeneous Equilibria, Magnitude of $K$, and $Q$}{Zumdahl \S13.4--13.5}

\begin{objectives}
  \item Write $K$ for heterogeneous equilibria.
  \item Interpret the magnitude of $K$, and use $Q$ to predict direction.
\end{objectives}

\reading{Zumdahl \S13.4--13.5}{660--671}

\begin{warmup}[4]
  \item $K$ expression for \ce{2 A <=> B}:
        \answerblank[3.4cm]{$[\ce{B}]/[\ce{A}]^2$}
  \item $\Delta n$ for \ce{2 NO2 <=> N2O4}: \answerblank[1.6cm]{$-1$}
  \item If $K = 100$, $K$ for the reverse is:
        \answerblank[1.8cm]{0.010}
  \item Do equilibrium constants carry units? \answerblank[1.4cm]{no}
\end{warmup}

\segment{INSTRUCTION A}{25}{Leaving things out}

\notesectionz{Heterogeneous equilibria}{13.4}
\practice{6}

When a reaction involves more than one phase, \term{pure solids} and
\term{pure liquids} are \answerblank[2cm]{omitted} from the $K$ expression.

The reason is not arbitrary: the ``concentration'' of a pure condensed phase
is fixed by its \answerblank[1.6cm]{density}, which does not change as the
reaction proceeds. Adding more solid does not make it more concentrated --- it
just makes the pile bigger.

\begin{workedexample}[Worked example: limestone decomposition]
\[ \ce{CaCO3(s) <=> CaO(s) + CO2(g)} \]
Both calcium-containing species are pure solids, so
\[ K_p = P_{\ce{CO2}} \]
The equilibrium pressure of \ce{CO2} above the solid depends only on
temperature. Doubling the amount of \ce{CaCO3} in the flask changes
\emph{nothing} --- provided some of each solid remains present.
\end{workedexample}

\segment{GUIDED PRACTICE}{15}{Write heterogeneous expressions}

\begin{enumerate}[leftmargin=*,itemsep=0.35em]
  \item \ce{2 H2O(l) <=> 2 H2(g) + O2(g)}:
        \answerblank[4.6cm]{$K = [\ce{H2}]^2[\ce{O2}]$}
  \item \ce{C(s) + CO2(g) <=> 2 CO(g)}:
        \answerblank[4.6cm]{$K = [\ce{CO}]^2/[\ce{CO2}]$}
  \item \ce{NH4Cl(s) <=> NH3(g) + HCl(g)}:
        \answerblank[4.6cm]{$K = [\ce{NH3}][\ce{HCl}]$}
\end{enumerate}

\segment{INSTRUCTION B}{20}{What $K$ tells you --- and what it does not}

\notesectionz{The extent of reaction}{13.5}
\practice{6}

\renewcommand{\arraystretch}{1.35}
\noindent\begin{tabular}{@{}p{3.0cm}p{4.4cm}p{5.6cm}@{}}
\toprule
\textbf{Magnitude} & \textbf{At equilibrium} & \textbf{Position}\\
\midrule
$K \gg 1$ & mostly \answerblank[1.8cm]{products} &
  lies far to the \answerblank[1.4cm]{right}\\
$K \approx 1$ & appreciable amounts of both & near the middle\\
$K \ll 1$ & mostly \answerblank[1.8cm]{reactants} &
  lies far to the \answerblank[1.2cm]{left}\\
\bottomrule
\end{tabular}

\begin{apctrap}
Zumdahl is explicit about this and it is a favourite exam trap:
\textbf{the size of $K$ says nothing about how fast equilibrium is reached.}
Extent is thermodynamics; speed is kinetics, governed by $E_a$. A reaction
with $K = 10^{20}$ can take centuries --- diamond turning to graphite is
exactly that case.
\end{apctrap}

\notesub{The reaction quotient}

$Q$ has the identical algebraic form as $K$, but uses whatever
concentrations you have \emph{right now}, equilibrium or not. Comparing
them predicts which way the system must move:

\renewcommand{\arraystretch}{1.3}
\noindent\begin{tabular}{@{}p{2.6cm}p{4.2cm}p{6.0cm}@{}}
\toprule
\textbf{Comparison} & \textbf{Meaning} & \textbf{Shift}\\
\midrule
$Q < K$ & too few products & toward \answerblank[2.0cm]{products} (right)\\
$Q = K$ & at equilibrium & \answerblank[1.6cm]{no shift}\\
$Q > K$ & too many products & toward \answerblank[1.8cm]{reactants} (left)\\
\bottomrule
\end{tabular}

\segment{APPLICATION}{20}{Using $Q$}

For \ce{H2(g) + I2(g) <=> 2 HI(g)}, $K = 50.0$ at a given temperature.

\begin{enumerate}[leftmargin=*,itemsep=0.4em]
  \item A flask contains all three at \SI{1.0}{\Molar}. Compute $Q$ and
        predict the direction of change. \workspace[2cm]
        \keyonly{\begin{apcanswer}$Q = (1.0)^2/[(1.0)(1.0)] = 1.0$.
        $Q < K$, so the reaction proceeds to the
        right.\end{apcanswer}}
  \item Another flask has $[\ce{HI}] = \SI{2.0}{\Molar}$,
        $[\ce{H2}] = [\ce{I2}] = \SI{0.010}{\Molar}$. Compute $Q$ and
        predict. \workspace[2cm]
        \keyonly{\begin{apcanswer}$Q = (2.0)^2/[(0.010)(0.010)] =
        4.0\times10^{4}$. $Q > K$, so the reaction proceeds to the
        left.\end{apcanswer}}
  \item Explain why $Q$ is useful even though it is not a new equation.
        \answerlines[2]{It converts an equilibrium constant into a
        prediction: comparing $Q$ with $K$ tells you which direction a
        non-equilibrium mixture will move without solving anything.}
\end{enumerate}

\begin{exitticket}
A reaction has $K = 3\times10^{15}$ yet no product forms in a laboratory
over several hours. Explain how both facts can be true.
\answerlines[2]{$K$ describes the equilibrium position, not the rate. A very
large activation energy can make the reaction imperceptibly slow even though
it is strongly product-favoured.}
\end{exitticket}

% =====================================================================
\blocklesson{Solving Equilibrium Problems}{Zumdahl \S13.6}

\begin{objectives}
  \item Build and complete an ICE table.
  \item Solve for equilibrium concentrations, including by approximation.
\end{objectives}

\reading{Zumdahl \S13.6}{671--676}

\begin{warmup}[3]
  \item If $Q < K$, the reaction shifts:
        \answerblank[2.2cm]{right}
  \item $K$ expression for \ce{CaCO3(s) <=> CaO(s) + CO2(g)}:
        \answerblank[3cm]{$K_p = P_{\ce{CO2}}$}
  \item Does a large $K$ mean a fast reaction? \answerblank[1.4cm]{no}
\end{warmup}

\segment{INSTRUCTION A}{25}{The ICE table}

\notesectionz{BCA tables, one chapter later}{13.6}
\practice{5}

The tool is the Chapter 3 BCA table with new labels:
\term{Initial}, \term{Change}, \term{Equilibrium}. As before, the change row
is set by the \answerblank[3.5cm]{coefficient ratios}; the only new feature
is that the reaction stops partway, so the change is an unknown $x$ rather
than a known quantity.

\renewcommand{\arraystretch}{1.45}
\noindent\begin{tabular}{@{}lccc@{}}
\toprule
 & \ce{N2O4} & $\rightleftharpoons$ & \ce{2 NO2}\\
\midrule
\textbf{Initial} & 1.00 & & 0\\
\textbf{Change} & \answerblank[1.4cm]{$-x$} & & \answerblank[1.4cm]{$+2x$}\\
\textbf{Equilibrium} & \answerblank[1.8cm]{$1.00-x$} & &
  \answerblank[1.4cm]{$2x$}\\
\bottomrule
\end{tabular}

\begin{workedexample}[Worked example 1: with the approximation]
\ce{N2O4(g) <=> 2 NO2(g)}, $K_c = 4.6\times10^{-3}$, starting at
$[\ce{N2O4}] = \SI{1.00}{\Molar}$.
\[
  K = \frac{(2x)^2}{1.00-x} = 4.6\times10^{-3}
\]
Because $K$ is small, very little reacts --- so assume
$1.00 - x \approx 1.00$:
\[
  4x^2 = 4.6\times10^{-3} \Rightarrow x^2 = 1.15\times10^{-3}
  \Rightarrow x = 0.0339
\]
\textbf{Validity check (the 5\% rule):} $0.0339/1.00 = 3.4\%$, comfortably
under 5\%, so the approximation stands.
\[
  [\ce{NO2}] = 2x = \SI{0.068}{\Molar}, \qquad
  [\ce{N2O4}] = \SI{0.966}{\Molar}
\]
\end{workedexample}

\begin{apctrap}
The approximation is only legal when $x$ is small compared with the initial
concentration --- roughly when $K$ is small. \textbf{Always run the 5\%
check.} If $x$ exceeds 5\% of the initial value, discard the approximation
and solve the quadratic properly.
\end{apctrap}

\segment{GUIDED PRACTICE}{15}{Set up the table}

For \ce{2 SO2 + O2 <=> 2 SO3} starting with \SI{2.0}{\Molar} \ce{SO2},
\SI{1.0}{\Molar} \ce{O2}, no \ce{SO3}, write the equilibrium row.
\answerlines[2]{$[\ce{SO2}] = 2.0-2x$; $[\ce{O2}] = 1.0-x$;
$[\ce{SO3}] = 2x$.}

\segment{INSTRUCTION B}{20}{When you must use the quadratic}

\notesectionz{Zumdahl's hydrogen fluoride problem}{13.6}
\practice{5}

\begin{workedexample}[Worked example 2: large $K$, no shortcut]
\ce{H2(g) + F2(g) <=> 2 HF(g)}, $K = 1.15\times10^{2}$, starting with
$[\ce{H2}] = \SI{1.000}{\Molar}$ and $[\ce{F2}] = \SI{2.000}{\Molar}$.

ICE gives $[\ce{H2}] = 1.000-x$, $[\ce{F2}] = 2.000-x$, $[\ce{HF}] = 2x$:
\[
  \frac{(2x)^2}{(1.000-x)(2.000-x)} = 115
\]
$K$ is large, so $x$ will \emph{not} be small --- the approximation is
unavailable. Expanding gives
\[
  111x^2 - 345x + 230 = 0
  \qquad\Rightarrow\qquad
  x = 0.968 \ \text{ or } \ 2.140
\]
\textbf{Reject $x = 2.140$}: subtracting it from 1.000 M would give a
negative concentration of \ce{H2}, which is physically impossible. So
\[
  [\ce{H2}] = \SI{0.032}{\Molar},\quad
  [\ce{F2}] = \SI{1.032}{\Molar},\quad
  [\ce{HF}] = \SI{1.936}{\Molar}
\]
\textbf{Reality check:} substituting back gives
$(1.936)^2/[(0.032)(1.032)] = 1.13\times10^{2}$, agreeing with the given
$1.15\times10^{2}$ to within rounding. \checkmark
\end{workedexample}

\begin{apcnote}
Two habits from that example are worth keeping. First, a quadratic always
gives two roots and one is normally \emph{physically} impossible --- reject
it by checking for negative concentrations, not by preference. Second,
always substitute your answer back into the $K$ expression. It takes ten
seconds and catches nearly every arithmetic slip.
\end{apcnote}

\segment{APPLICATION}{20}{Full problem}

\ce{H2(g) + I2(g) <=> 2 HI(g)} has $K = 50.0$. A flask initially contains
\SI{1.00}{\Molar} \ce{H2} and \SI{1.00}{\Molar} \ce{I2}.

\begin{enumerate}[leftmargin=*,itemsep=0.4em]
  \item Build the ICE table.
        \answerlines[2]{$[\ce{H2}] = [\ce{I2}] = 1.00-x$;
        $[\ce{HI}] = 2x$.}
  \item Solve for $x$. (This one is a perfect square --- take the square
        root of both sides.) \workspace[2.6cm]
        \keyonly{\begin{apcanswer}$(2x)^2/(1.00-x)^2 = 50.0$;
        $2x/(1.00-x) = 7.071$; $2x = 7.071 - 7.071x$;
        $9.071x = 7.071$; $x = 0.779$\end{apcanswer}}
  \item State all three equilibrium concentrations and check them against
        $K$. \workspace[2.2cm]
        \keyonly{\begin{apcanswer}$[\ce{H2}] = [\ce{I2}] =
        \SI{0.221}{\Molar}$; $[\ce{HI}] = \SI{1.558}{\Molar}$.
        Check: $(1.558)^2/(0.221)^2 = 49.7 \approx 50.0$
        \checkmark\end{apcanswer}}
\end{enumerate}

\begin{exitticket}
When is it safe to assume $x$ is negligible, and how do you confirm it?
\answerlines[2]{When $K$ is small, so very little reaction occurs. Confirm
by checking that $x$ is less than 5\% of the initial concentration it was
subtracted from.}
\end{exitticket}

% =====================================================================
\blocklesson{Le Ch\^{a}telier's Principle}{Zumdahl \S13.7}

\begin{objectives}
  \item Predict the direction of shift for concentration, volume, and
        temperature changes.
  \item Explain which changes alter $K$ and which do not.
\end{objectives}

\reading{Zumdahl \S13.7}{676--683}

\begin{warmup}[3]
  \item The 5\% rule checks what? \answerblank[6.0cm]{validity of the
        approximation}
  \item Reject a quadratic root when it gives:
        \answerblank[5.0cm]{a negative concentration}
  \item If $Q > K$, the shift is: \answerblank[2cm]{left}
\end{warmup}

\segment{INSTRUCTION A}{25}{The principle}

\notesectionz{Systems relieve stress}{13.7}
\practice{6}

\begin{apcnote}
\textbf{Le Ch\^{a}telier's principle:} if a change is imposed on a system at
equilibrium, the equilibrium position shifts in the direction that
partially \answerblank[2cm]{offsets} that change.

Zumdahl's formulation for concentration is particularly clean: add a
component and the system shifts in the direction that
\emph{lowers} the concentration of that component. Remove one and the
opposite happens.
\end{apcnote}

\renewcommand{\arraystretch}{1.45}
\noindent\begin{tabular}{@{}p{3.6cm}p{4.4cm}p{4.8cm}@{}}
\toprule
\textbf{Change} & \textbf{Shift} & \textbf{Does $K$ change?}\\
\midrule
Add a reactant & toward \answerblank[2.0cm]{products} &
  \answerblank[1.2cm]{no}\\
Remove a product & toward \answerblank[2.0cm]{products} &
  \answerblank[1.2cm]{no}\\
Decrease volume & toward the side with \answerblank[3.3cm]{fewer gas moles} &
  \answerblank[1.2cm]{no}\\
Increase temperature & toward the \answerblank[2.6cm]{endothermic} direction &
  \textbf{\answerblank[1.2cm]{yes}}\\
Add a catalyst & \answerblank[1.8cm]{no shift} &
  \answerblank[1.2cm]{no}\\
Add an inert gas at constant $V$ & \answerblank[1.8cm]{no shift} &
  \answerblank[1.2cm]{no}\\
\bottomrule
\end{tabular}

\begin{apctrap}
\textbf{Only temperature changes the value of $K$.} Everything else shifts
the \emph{position} of equilibrium while leaving the constant alone --- the
system simply moves to restore the same $K$.

Two more traps in that table. A catalyst speeds both directions equally, so
equilibrium arrives sooner but sits in the same place. And an inert gas
added at constant volume changes the total pressure without changing any
\emph{partial} pressure, so nothing shifts.
\end{apctrap}

\segment{GUIDED PRACTICE}{15}{Predict the shifts}

For \ce{N2(g) + 3 H2(g) <=> 2 NH3(g)}, $\Delta H = \SI{-92}{\kilo\joule}$:
\begin{enumerate}[leftmargin=*,itemsep=0.3em]
  \item Add \ce{N2}: \answerblank[2.4cm]{right}
  \item Remove \ce{NH3}: \answerblank[2.4cm]{right}
  \item Decrease the volume: \answerblank[4.6cm]{right (4 mol $\to$ 2 mol)}
  \item Increase the temperature:
        \answerblank[5.2cm]{left (exothermic forward)}
  \item Add a catalyst: \answerblank[2.4cm]{no shift}
\end{enumerate}

\segment{INSTRUCTION B}{20}{Temperature and heterogeneous cases}

\notesectionz{Two special situations}{13.7}
\practice{4}

\notesub{Treat heat as a term in the equation}

For an exothermic reaction, write heat as a \answerblank[1.6cm]{product}:
\[ \ce{N2 + 3 H2 <=> 2 NH3} + \text{heat} \]
Now Le Ch\^{a}telier applies directly --- adding heat (raising the
temperature) pushes the system \answerblank[1.4cm]{left}, exactly as adding
any other product would.

\begin{workedexample}[Worked example: Zumdahl's arsenic roasting]
\[ \ce{As4O6(s) + 6 C(s) <=> As4(g) + 6 CO(g)} \]
\textbf{(a) Add \ce{CO}:} shift \emph{left}, away from the added substance.

\textbf{(b) Add or remove \ce{C} or \ce{As4O6}:} \emph{no shift at all}.
Both are pure solids, so they do not appear in $K$ --- adding more changes
no concentration. (Removing \emph{all} of one would stop the reaction, but
that is a different situation.)

\textbf{(c) Remove \ce{As4} gas:} shift \emph{right}, to replace it.
\end{workedexample}

\segment{APPLICATION}{20}{Applied Le Ch\^{a}telier}

\begin{enumerate}[leftmargin=*,itemsep=0.4em]
  \item For \ce{2 SO2(g) + O2(g) <=> 2 SO3(g)}, $\Delta H < 0$: state one
        change that increases the yield of \ce{SO3} and one that does not,
        with reasons. \answerlines[3]{Increases yield: raise the pressure
        (fewer gas moles on the right), or lower the temperature (shifts
        toward the exothermic direction). Does not: add a catalyst --- it
        speeds both directions and leaves the position unchanged.}
  \item The Haber process is run at high temperature even though that
        \emph{lowers} the equilibrium yield of ammonia. Explain the
        trade-off. \answerlines[3]{Low temperature favours the position of
        equilibrium but makes the reaction far too slow to be useful.
        Industry accepts a smaller equilibrium yield in exchange for
        reaching it quickly --- a kinetics-versus-thermodynamics compromise,
        partly recovered by using high pressure and a catalyst.}
  \item Explain why adding argon to a rigid vessel containing an
        equilibrium mixture causes no shift.
        \answerlines[2]{At constant volume the partial pressures of the
        reacting gases are unchanged, so $Q$ still equals $K$ --- only the
        total pressure rises.}
\end{enumerate}

\begin{exitticket}
Which single type of change alters the numerical value of $K$, and why do
the others not? \answerlines[2]{Temperature. All other changes disturb the
concentrations, and the system shifts precisely to restore the same $K$;
temperature changes the constant itself because it changes the underlying
thermodynamics.}
\end{exitticket}

\end{document}
